% Successor v10, review-only M2 verification note, 11 October 2026.
% This file is not automatically inserted into main.tex.
% The four author-controlled repairs and all mathematical prose are frozen.
% PENDING exact-head Lean and axiom audit; do not mark GREEN yet.
%
% Candidate insertion at Lemma \ref{lem:level-removal} (label must be checked
% against the actual author v10 source before applying any patch):
\todo[inline]{\textbf{v10 validation: level removal, complete type transport
-- PENDING.}
The existing finite \texttt{LevelRemovalPartial.lean} and
\texttt{LevelRemovalFreeCut.lean} establish the strict spacing guard
$\mathrm{fl}_A(m+1)<m$, preservation of the forbidden $L$-age, and
the exact new free cut after deleting ordinary coordinate $m$.
The proposed
\texttt{deleteAtPartial\_partialTypeAt\_eq\_deleteOrdinary}
in \texttt{LevelRemovalType.lean} additionally identifies the
\emph{complete} $L^+$ type of each retained vertex with deletion
of coordinate $m$ from its previous complete type, retaining
both orientations and absences of every binary tuple, all
unary/diagonal facts, and all $E$ pairs. Its proof is a direct
induced-model calculation and imposes no new age hypothesis.

This remains a \emph{finite local} lemma. To obtain
Lemma~\ref{lem:level-removal} globally, one must prove that the
deletion position $m=\widetilde F(n)-1$ is guarded in every relevant
image $F(T)$, that the resulting image types are admissible
at their transported free cuts, and that the uniform map preserves
predecessors and weak successors. For M2 one must then verify
the corrected finite $L$-only B3 for the induced prescription;
the earlier $E$-partial image-transport statement alone is
insufficient. Do not mark M2 or Proposition~\ref{prob:upperbound}
GREEN from this lemma.}
